\documentclass[11pt,aspectratio=169]{beamer}

\usetheme{Madrid}
\usecolortheme{seahorse}

\usepackage[utf8]{inputenc}
\usepackage{amsmath,amssymb}
\usepackage{algpseudocode}
\usepackage{algorithmicx}
\usepackage{xcolor}
\usepackage{listings}
\usepackage{booktabs}
\usepackage{tikz}
\usetikzlibrary{arrows.meta,positioning,fit,backgrounds}

\setbeamertemplate{navigation symbols}{}
\setbeamertemplate{footline}[frame number]

\lstdefinestyle{py}{
  language=Python,
  basicstyle=\ttfamily\scriptsize,
  keywordstyle=\color{blue!70!black}\bfseries,
  commentstyle=\color{gray},
  stringstyle=\color{green!40!black},
  showstringspaces=false,
  breaklines=true,
  numbers=left,
  numberstyle=\tiny\color{gray},
  frame=single,
  xleftmargin=1.5em,
  aboveskip=0.3em,
  belowskip=0.3em
}

\title{LeetCode Deep Dive}
\subtitle{From First Instinct to Optimal Algorithm: Dynamic Programming in Practice\\
\small Week 8 Supplement -- TOAE209/TOAH209: Design and Analysis of Algorithms}
\author{Foreign Trade University}
\date{}

\begin{document}

% ---------------------------------------------------------------
\begin{frame}
  \titlepage
\end{frame}
% ---------------------------------------------------------------

\begin{frame}{Outline}
  \tableofcontents
\end{frame}

% =================================================================
\section{A Framework for Attacking a New DP Problem}
% =================================================================

\begin{frame}{Why walk through LeetCode problems?}
  \begin{itemize}
    \item This week's \texttt{LEETCODE.md} list gives ten Week 8 problems to practice
    on your own. This deck narrates the thinking process for two of them, end to end,
    so you can reuse it on the rest.
    \item We pick one \textcolor{blue!50!black}{\textbf{Medium}} (\emph{Coin Change} --
    unbounded knapsack) and one \textcolor{red!60!black}{\textbf{Hard}} (\emph{Maximum
    Profit in Job Scheduling} -- weighted interval scheduling, this week's \emph{other}
    lecture topic, revisited).
    \item Both times the naive instinct is the same one the lecture already warned
    about: a recursive solution with \textbf{overlapping subproblems} that recomputes
    the same answer over and over -- exactly the Fibonacci blow-up from this week's
    notes.
  \end{itemize}
\end{frame}

\begin{frame}{The four-step process we will repeat twice}
  \begin{enumerate}
    \item \textbf{First instinct.} Write the recursive recurrence that follows directly
    from the problem statement, with no memoization.
    \item \textbf{Measure why it hurts.} Count how many times the same subproblem gets
    recomputed, and see the exponential blow-up.
    \item \textbf{Find the structural insight.} Identify the \emph{subproblem} (this
    week's step 1 of the 5-step DP recipe) whose optimal-substructure recurrence
    (step 2) lets one pass fill every value exactly once.
    \item \textbf{Implement the optimal algorithm}, trace the DP table by hand on a
    small example, then look for the programming-level tricks that make the
    implementation itself correct and fast.
  \end{enumerate}
  \begin{alertblock}{This mirrors the course}
    Steps 1--3 are exactly this week's ``define the subproblem, write the recurrence,
    analyze the cost'' recipe; step 4 is what turns a correct recurrence into working
    code.
  \end{alertblock}
\end{frame}

% =================================================================
\section{Problem 1 (Medium): Coin Change}
% =================================================================

\begin{frame}{Problem statement}
  \begin{block}{LeetCode 322 -- Coin Change (\textcolor{blue!50!black}{Medium})}
    You are given coin denominations \texttt{coins} (an unlimited supply of each) and a
    target \texttt{amount}. Return the \textbf{fewest} coins needed to make up
    \texttt{amount}, or $-1$ if it cannot be made up at all.
  \end{block}
  \begin{itemize}
    \item Constraints (LeetCode): $1 \le \texttt{coins.length} \le 12$,
    $1 \le \texttt{amount} \le 10^4$.
    \item This is exactly the \textbf{unbounded knapsack} variant this week's lecture
    covers under ``Coin change: minimum number of coins.''
  \end{itemize}
\end{frame}

\begin{frame}{Worked example}
  \begin{block}{\texttt{coins = [1, 2, 5]}, \texttt{amount = 11}}
    \textbf{Answer: 3} -- e.g.\ $5+5+1 = 11$ using three coins. (Also $5+2+2+2=11$ uses
    four -- more, so it's not optimal.)
  \end{block}
  \begin{itemize}
    \item Greedy ``always take the largest coin that fits'' happens to work for US
    coinage but is not correct in general (e.g.\ \texttt{coins=[1,3,4]},
    \texttt{amount=6}: greedy gives $4+1+1=3$ coins, but $3+3=2$ coins is better). We
    need an algorithm that is \emph{provably} optimal for \emph{any} coin set.
  \end{itemize}
\end{frame}

\begin{frame}{First instinct: recurse on ``which coin last?''}
  \begin{block}{Most direct reading of the problem}
    $\textsc{MinCoins}(a) = 1 + \min_{c \,\in\, \texttt{coins},\, c \le a}
    \textsc{MinCoins}(a-c)$, with $\textsc{MinCoins}(0)=0$.
  \end{block}
  \begin{algorithmic}[1]
    \Function{MinCoins}{$a$}
      \If{$a = 0$} \Return $0$ \EndIf
      \If{$a < 0$} \Return $\infty$ \EndIf
      \State \Return $1 + \min_{c \in \texttt{coins}} \Call{MinCoins}{a - c}$
    \EndFunction
  \end{algorithmic}
  This is a direct transcription of the problem -- and it is \textbf{correct}.
\end{frame}

\begin{frame}{Why this is bad}
  \begin{itemize}
    \item \texttt{MinCoins(11)} calls \texttt{MinCoins(10)}, \texttt{MinCoins(9)},
    \texttt{MinCoins(6)}; \texttt{MinCoins(10)} \emph{also} calls \texttt{MinCoins(9)}
    and \texttt{MinCoins(8)} and \texttt{MinCoins(5)} -- the same smaller amounts get
    recomputed from many different paths.
    \item This is \textbf{exactly} the overlapping-subproblems blow-up this week's
    notes illustrate with Fibonacci ($F(n)$ makes $2F(n+1)-1$ calls without
    memoization). Here the branching factor is $|\texttt{coins}|$ and the depth is up
    to \texttt{amount}, so the call count grows like $|\texttt{coins}|^{\texttt{amount}}$
    in the worst case.
    \item At \texttt{amount} $=10^4$ this is not a slowdown -- it is astronomically
    infeasible.
  \end{itemize}
  \begin{alertblock}{The smell to notice}
    Whenever a recursion's arguments repeat across different call paths, that is the
    signature of overlapping subproblems -- memoize by that argument.
  \end{alertblock}
\end{frame}

\begin{frame}{The structural insight}
  \begin{itemize}
    \item There are only \texttt{amount}$+1$ \emph{distinct} subproblems
    ($a = 0, 1, \dots, \texttt{amount}$), each doing $O(|\texttt{coins}|)$ work once its
    smaller answers are known.
    \item \textbf{Subproblem (step 1 of this week's recipe):} $dp[a]$ = fewest coins to
    make amount $a$.
    \item \textbf{Recurrence (step 2):} $dp[a] = 1 + \min_{c \le a} dp[a-c]$,
    $dp[0]=0$, evaluated bottom-up for $a = 1, \dots, \texttt{amount}$ so every smaller
    value is already known -- no recomputation, ever.
  \end{itemize}
\end{frame}

\begin{frame}{Trace on the worked example}
  \scriptsize
  \texttt{coins = [1, 2, 5]} -- $dp[a]$ for $a = 0, \dots, 11$ (bottom-up):
  \begin{center}
  \begin{tabular}{@{}c*{12}{c}@{}}
    \toprule
    $a$ & 0 & 1 & 2 & 3 & 4 & 5 & 6 & 7 & 8 & 9 & 10 & 11 \\
    \midrule
    $dp[a]$ & 0 & 1 & 1 & 2 & 2 & 1 & 2 & 2 & 3 & 3 & 2 & \textbf{3} \\
    \bottomrule
  \end{tabular}
  \end{center}
  \begin{itemize}
    \item $dp[5]=1$ (one 5-coin) is what makes $dp[10]=2$ ($5+5$) and, in turn,
    $dp[11] = 1+dp[10] = 1 + 2 = 3$ ($5+5+1$) beat $1+dp[9]=4$ and $1+dp[6]=3$ (tied,
    both valid optimal decompositions).
    \item Result: $dp[11] = 3$, matching the expected answer.
  \end{itemize}
\end{frame}

\begin{frame}{Complexity: naive vs.\ optimal}
  \begin{center}
  \begin{tabular}{@{}lll@{}}
    \toprule
    Approach & Time & At \texttt{amount}$=10^4$ \\
    \midrule
    Unmemoized recursion & $O(|\texttt{coins}|^{\texttt{amount}})$ & astronomically
    infeasible \\
    \textbf{Bottom-up DP} & $\boldsymbol{O(\texttt{amount}\times|\texttt{coins}|)}$ &
    $\le 1.2\times10^5$ ops \\
    \bottomrule
  \end{tabular}
  \end{center}
  Exactly the same shape of improvement as this week's Knapsack DP: exponential
  recursion collapses to (number of subproblems) $\times$ (work per subproblem) once
  each subproblem is solved only once.
\end{frame}

\begin{frame}[fragile]{Python implementation}
  \begin{lstlisting}[style=py]
def coinChange(coins, amount):
    INF = amount + 1                  # sentinel: "not (yet) reachable"
    dp = [0] + [INF] * amount         # dp[a] = fewest coins for amount a

    for a in range(1, amount + 1):
        for c in coins:
            if c <= a:
                dp[a] = min(dp[a], dp[a - c] + 1)

    return dp[amount] if dp[amount] != INF else -1
  \end{lstlisting}
\end{frame}

\begin{frame}{Implementation tips \& tricks (the programming side)}
  \small
  \begin{itemize}
    \item \textbf{Sentinel choice: \texttt{amount + 1}, not \texttt{float('inf')}.}
    \texttt{amount+1} is already larger than any valid answer (which is at most
    \texttt{amount}, using all 1-coins), and unlike \texttt{float('inf')} it stays an
    \texttt{int}, so \texttt{dp[a-c] + 1} never silently becomes a float.
    \item \textbf{Iterative DP, not memoized recursion, for the same reason as
    Week 3's Union-Find.} A recursive \texttt{@lru\_cache} version recurses to depth
    \texttt{amount} in the worst case (e.g.\ \texttt{coins=[1]}); at
    \texttt{amount}$=10^4$ that exceeds Python's default recursion limit
    ($\approx 1000$) and crashes. The bottom-up loop has no such limit.
    \item \textbf{Loop order doesn't matter here, but it would for ``Coin Change II.''}
    Because each coin can be reused arbitrarily (unbounded knapsack), iterating
    \texttt{amount} outside and \texttt{coins} inside is safe. In the closely related
    \emph{counting} problem (number of \emph{combinations}, not permutations), swapping
    the loop order changes the answer -- a classic, easy-to-miss bug.
  \end{itemize}
\end{frame}

% =================================================================
\section{Problem 2 (Hard): Maximum Profit in Job Scheduling}
% =================================================================

\begin{frame}{Problem statement}
  \begin{block}{LeetCode 1235 -- Maximum Profit in Job Scheduling
  (\textcolor{red!60!black}{Hard})}
    You are given \texttt{n} jobs, each with a \texttt{startTime}, \texttt{endTime}, and
    \texttt{profit}. You may schedule any subset of \textbf{non-overlapping} jobs
    (job $A$ and job $B$ don't overlap if $A$'s end $\le$ $B$'s start, or vice versa).
    Return the \textbf{maximum total profit}.
  \end{block}
  \begin{itemize}
    \item Constraints (LeetCode): $1 \le n \le 5\times10^4$.
    \item This is \textbf{exactly} this week's other lecture topic -- Weighted Interval
    Scheduling -- with \texttt{profit} playing the role of $w_j$.
  \end{itemize}
\end{frame}

\begin{frame}{Worked example}
  \begin{block}{\texttt{startTime=[1,2,3,3]}, \texttt{endTime=[3,4,5,6]},
  \texttt{profit=[50,10,40,70]}}
    \textbf{Answer: 120} -- take job $1$ ($[1,3)$, profit $50$) and job $4$ ($[3,6)$,
    profit $70$): they don't overlap ($3 \le 3$), total $50+70=120$. Any other
    combination scores less (e.g.\ job $3$ alone plus job $4$ overlaps, since $3 < 5$
    is not $\le 3$).
  \end{block}
\end{frame}

\begin{frame}{First instinct, and why the naive version is bad}
  \small
  \begin{block}{Most direct reading of the problem}
    Try every subset of jobs, keep the non-overlapping ones with maximum profit sum.
  \end{block}
  \begin{itemize}
    \item That is $O(2^n)$ -- utterly infeasible at $n = 5\times10^4$.
    \item Even the ``obviously smarter'' fix -- sort by end time, then for each job
    scan \emph{all previous} jobs to find the latest one that ends before it starts --
    is a real DP (correct!), but costs $O(n)$ per job for that scan, so
    $O(n^2)$ total. At $n=5\times10^4$ that's $2.5\times10^9$ operations: still too
    slow.
  \end{itemize}
  \begin{alertblock}{The smell to notice}
    A DP whose recurrence is right but whose \emph{per-subproblem} work is a linear
    scan is a sign that the scan itself can probably be replaced by a binary search --
    exactly this week's lecture on Weighted Interval Scheduling.
  \end{alertblock}
\end{frame}

\begin{frame}{The structural insight}
  \begin{itemize}
    \item Sort jobs by \textbf{end time}. Define $p(j)$ = the largest index $i < j$
    with $\text{end}_i \le \text{start}_j$ (the latest job compatible with $j$).
    \item \textbf{Subproblem:} $M[j]$ = max profit using only jobs $1,\dots,j$.
    \item \textbf{Recurrence:} $M[j] = \max\big(M[j-1],\ \text{profit}_j + M[p(j)]\big)$
    -- either skip job $j$, or take it and add the best profit compatible with it.
    \item Because jobs are sorted by end time, $p(j)$ is exactly ``the rightmost
    already-processed end time that is $\le \text{start}_j$'' -- a \textbf{binary
    search} into the (already sorted) array of end times, in $O(\log n)$.
  \end{itemize}
\end{frame}

\begin{frame}{Trace on the worked example}
  \small
  Sorted by end time: job1$(1,3,50)$, job2$(2,4,10)$, job3$(3,5,40)$, job4$(3,6,70)$.
  \begin{center}
  \begin{tabular}{@{}cccccl@{}}
    \toprule
    $j$ & start$_j$ & profit$_j$ & $p(j)$ & $M[j]$ & reasoning \\
    \midrule
    1 & 1 & 50 & 0 & $50$ & $\max(M[0]{=}0,\ 50+M[0]{=}50)$ \\
    2 & 2 & 10 & 0 & $50$ & $\max(M[1]{=}50,\ 10+M[0]{=}10)$ \\
    3 & 3 & 40 & 1 & $90$ & $\max(M[2]{=}50,\ 40+M[1]{=}90)$ \\
    4 & 3 & 70 & 1 & \textbf{120} & $\max(M[3]{=}90,\ 70+M[1]{=}120)$ \\
    \bottomrule
  \end{tabular}
  \end{center}
  $p(3)=p(4)=1$ because job1 is the latest job ending at or before start $3$.
  Result: $M[4] = 120$, matching the expected answer.
\end{frame}

\begin{frame}{Complexity: naive vs.\ optimal}
  \begin{center}
  \begin{tabular}{@{}lll@{}}
    \toprule
    Approach & Time & At $n=5\times10^4$ \\
    \midrule
    Try all subsets & $O(2^n)$ & infeasible \\
    DP with linear scan for $p(j)$ & $O(n^2)$ & $\sim 2.5\times10^9$ ops \\
    \textbf{DP with binary search for $p(j)$} & $\boldsymbol{O(n\log n)}$ &
    $\sim 8\times10^5$ ops \\
    \bottomrule
  \end{tabular}
  \end{center}
  Same $O(n\log n)$ bound as this week's lecture proves for Weighted Interval
  Scheduling in general.
\end{frame}

\begin{frame}[fragile]{Python implementation}
  \begin{lstlisting}[style=py]
from bisect import bisect_right

def jobScheduling(startTime, endTime, profit):
    jobs = sorted(zip(endTime, startTime, profit))   # sort by end time
    ends = [e for e, s, p in jobs]                   # already sorted, reused below

    n = len(jobs)
    M = [0] * (n + 1)                # M[j] = best profit using jobs[0:j]
    for j in range(1, n + 1):
        end_j, start_j, profit_j = jobs[j - 1]
        i = bisect_right(ends, start_j, 0, j - 1)     # p(j): jobs[0:i] all end <= start_j
        M[j] = max(M[j - 1], profit_j + M[i])

    return M[n]
  \end{lstlisting}
\end{frame}

\begin{frame}{Implementation tips \& tricks (the programming side)}
  \small
  \begin{itemize}
    \item \textbf{Sort the three parallel arrays \emph{together}.} A common bug is
    sorting \texttt{endTime} alone and forgetting to permute \texttt{startTime} and
    \texttt{profit} to match. Zipping all three into tuples before sorting makes this
    impossible to get wrong.
    \item \textbf{Reuse \texttt{bisect\_right} instead of hand-rolling binary search.}
    \texttt{bisect\_right(ends, start\_j, 0, j-1)} directly returns the count of already
    -processed jobs whose end time is $\le \text{start}_j$ -- which \emph{is} $p(j)$ in
    this 1-indexed $M$ array, with zero off-by-one arithmetic to get wrong.
    \item \textbf{The bound \texttt{0, j-1} matters.} Searching only the \emph{prefix}
    already sorted by end time (jobs $1,\dots,j-1$) is both correct and avoids
    accidentally matching a later job that happens to share the same end time as job
    $j$ itself.
  \end{itemize}
\end{frame}

% =================================================================
\section{Summary}
% =================================================================

\begin{frame}{Summary: two problems, one lesson}
  \footnotesize
  \begin{enumerate}
    \item \textbf{Coin Change} is unbounded knapsack: the naive recursion recomputes
    the same amount from many paths (this week's Fibonacci warning, restated), and
    bottom-up DP fixes it in $O(\texttt{amount}\times|\texttt{coins}|)$.
    \item \textbf{Maximum Profit in Job Scheduling} is Weighted Interval Scheduling with
    a different name: the DP recurrence is already optimal, but a linear scan for
    $p(j)$ silently costs $O(n^2)$ -- replacing it with binary search recovers the
    lecture's $O(n\log n)$ bound.
    \item Both times, the naive instinct was correct but wasteful, in the \emph{exact}
    shape this week's 5-step DP recipe predicts: define the subproblem, write the
    recurrence, and only then does the running time become (subproblems) $\times$
    (work per subproblem) -- and that per-subproblem work is itself worth optimizing.
    \item The implementation tricks (integer sentinels, iterative not recursive DP,
    \texttt{bisect} instead of hand-rolled search) are a second, independent layer:
    even an asymptotically optimal recurrence can fail in practice if the code hits
    language-level limits or off-by-one errors.
  \end{enumerate}
  \begin{alertblock}{Keep practicing}
    \texttt{LEETCODE.md} lists eight more Week 8 problems (\emph{House Robber},
    \emph{Partition Equal Subset Sum}, \emph{Target Sum}, \emph{Ones and Zeroes},
    \dots) -- try applying the same four-step process to each of them.
  \end{alertblock}
\end{frame}

\end{document}
